\begin{helper}
Let \((\Omega,\nu)\) be a measure space, and let \(F,C,E\subseteq\Omega\).
Assume \(E\subseteq F\), \(\nu(F)<\infty\), and \(\nu(F\cap C)=\nu(F)\).
If \(\nu(E\cap C)\le B\), then \(\nu(E)\le B\).
\end{helper}
\begin{proof}
Since \(F\cap C\subseteq F\) and the two sets have the same finite measure,
the part of \(F\) outside \(C\) has measure zero.  More explicitly, by finite
additivity on the disjoint decomposition
\[
F=(F\cap C)\,\dot\cup\, (F\setminus C)
\]
and by the hypothesis \(\nu(F)<\infty\), subtracting the equality
\(\nu(F\cap C)=\nu(F)\) gives \(\nu(F\setminus C)=0\).  Because
\(E\subseteq F\), the outside-cube part \(E\setminus C\) is contained in
\(F\setminus C\), and hence \(\nu(E\setminus C)=0\).  Applying finite
additivity to
\[
E=(E\cap C)\,\dot\cup\, (E\setminus C)
\]
therefore gives \(\nu(E)=\nu(E\cap C)\).  The assumed bound
\(\nu(E\cap C)\le B\) now implies \(\nu(E)\le B\).
\end{proof}
